\documentclass[11pt]{article}
\usepackage{fontspec,amsmath,amssymb,amsthm,hyperref}
\usepackage[margin=1in]{geometry}
\setlength{\emergencystretch}{3em}
\title{Canonical Gaussian law of the actual splitting algorithm}
\author{LangevinLean mathematical exposition}
\date{}
\begin{document}\maketitle
\noindent Source module: \texttt{ProximalLangevin/SplittingGaussianLaw.lean}. Compile this proof document with XeLaTeX or LuaLaTeX. The accompanying HTML places the complete Lean declarations beside every proof.


The scalar-coordinate Gaussian hypotheses determine the complete vector innovation law and its infinite iid path. Hence every realization with a fixed initial law gives the same actual PIPGLA output law.


\section{Assemble scalar coordinates into iid standard Gaussian vectors}


The coordinate laws and their joint independence imply \(\xi_k\sim\mathcal N(0,I)\) as a full state vector, and independence of these vectors across \(k\). Together with initialization/path independence, the generic iid transfer theorem gives the chain law driven by \(\mathcal N(0,I)^{\otimes\mathbb N}\).



\noindent Lean declarations: \texttt{pipglaLaws\_eq\_canonicalGaussian}.

\section{Use either canonical noise-path representation}


The countable product of scalar Gaussian coordinates, grouped by iteration, has the same noise-path pushforward law as the vector product. Apply the measurable chain to this identity. Two probability-space realizations with the same initialization law therefore both equal the same canonical Gaussian chain law at time \(n\), proving equality of their actual output laws.



\noindent Lean declarations: \texttt{pipglaLaws\_eq\_canonicalNoisePath}, \texttt{pipglaLaws\_eq\_of\_initialLaw\_eq}.

\end{document}
